%***************************************************************************************************

% Basic language settings used to format the document
%\usepackage[english,polish]{babel}
% \usepackage[english]{babel}
\usepackage[english]{babel}

% The babel package for Polish causes a conflict with the amssymb package.
% The '\lll' command is defined by both packages - the second definition is required.
% \let\lll\undefined

% For multilingual documents, sets the main document language
% \selectlanguage{english}

% Document encoding
\usepackage[utf8]{inputenc} 
% \usepackage[T1]{fontenc} 

% Any font size, character encoding
% \usepackage{lmodern}

% Paragraph indentation
\usepackage{indentfirst}

% Typographic formatting
\frenchspacing

% Provides numerous improvements for displaying and organising mathematical formulas.
\usepackage{amsmath}

% Introduces an extended set of symbols, including \leadsto
\usepackage{amssymb}

% Additional "script" mathematical font
\usepackage{mathrsfs}

% Additional support for the mathbb environment, which does not support digits by default (\mathbb{})
%\usepackage{bbold}

% Fixes/improves amsmath
\usepackage{mathtools}

\usepackage{graphicx}
\usepackage{subcaption}


% ***************************************************************************************************
% Colours
% ***************************************************************************************************

% Enables colouring of individual table cells
\usepackage[table, svgnames]{xcolor}% http://ctan.org/pkg/

% Enables easy changing of table line colours
\usepackage{tabu}

\usepackage{longtable}

% Provides extended control over colours.
\usepackage{xcolor}

% Colour definitions
\definecolor{lgray}{HTML}{9F9F9F}
\definecolor{dgray}{HTML}{5F5F5F}

% ***************************************************************************************************
% Algorithms
% ***************************************************************************************************

% Provides an environment for constructing pseudocode
\usepackage[ruled,vlined,linesnumbered,longend,algochapter]{algorithm2e}
\renewcommand{\algorithmcfname}{Algorithm}

% Changes the environment name from the default "Algorithm X" to "Pseudocode X"
\newenvironment{algorithm-custom}[1][htb]{ 
	% \renewcommand{\algorithmcfname}{Algorithm}
	\begin{algorithm}[#1]%
	}{
\end{algorithm}
}

% Changes the comment size
\newcommand\algcomment[1]{
	\footnotesize{#1}
}

% Sets the specified style for comments
\SetCommentSty{algcomment}

% Centred tilde
\usepackage{textcomp}%
\newcommand{\textapprox}{\raisebox{0.5ex}{\texttildelow}}

% Source code listing
\usepackage{listings}
\usepackage{xcolor}

% Style definition for Python code
\lstdefinestyle{pythonstyle}{
    language=Python,
    backgroundcolor=\color{gray!10},   % light grey background
    basicstyle=\ttfamily\small,        % monospaced font, small size
    keywordstyle=\color{blue}\bfseries, % keywords in blue
    commentstyle=\color{green!60!black}\itshape, % comments in green italics
    stringstyle=\color{red!70!black},  % strings in red
    numbers=left,                      % line numbering on the left
    numberstyle=\tiny\color{gray},     % line number style
    stepnumber=1,                       % number every line
    numbersep=8pt,                      % distance between numbers and code
    showstringspaces=false,             % do not show spaces in strings
    tabsize=4,                          % tab = 4 spaces
    breaklines=true,                    % wrap long lines
    breakatwhitespace=true,             % break lines at whitespace
    frame=single,                        % frame around the code
    rulecolor=\color{black},            % frame colour
    captionpos=b                        % caption below the code
}
\renewcommand{\lstlistingname}{Source code} % Listing name
% \renewcommand{\lstlistingname}{Source code} % Listing name

% ***************************************************************************************************
% Margins
% ***************************************************************************************************

% Page size and margin settings
\usepackage[headheight=18pt, top=25mm, bottom=25mm, left=25mm, right=25mm]{geometry}

% Removes the top margin for environments
\makeatletter
\setlength\@fptop{0\p@}	
\makeatother

% ***************************************************************************************************
% Style
% ***************************************************************************************************

% Defines the 'titlingpage' environment, which provides full control over the title page layout.
\usepackage{titling}

% Enables modification of the table of contents style
% \usepackage[subfigure]{tocloft}	
\usepackage{tocloft}	

\tocloftpagestyle{tableOfContentStyle}

% Defining custom header and/or footer styles
\usepackage{fancyhdr}

% Default style for the document
\fancypagestyle{custom}{
	\fancyhf{}									% clear footers and headers
	\fancyhead[RO]{								% Right, odd header
		\hrulefill \hspace{16pt} \large Chapter \thechapter
        %\hrulefill \hspace{16pt} \large Chapter \thechapter
		\put(-471.5,5.5){%
			\makebox(0,0)[l]{%
                \small Wrocław University of Science and Technology
				%\small Wrocław University of Science and Technology
			}
		}
	}
	\fancyhead[LE]{								% Left, even header
        \large Chapter \thechapter \hspace{16pt} \hrulefill 
		%\large Chapter \thechapter \hspace{16pt} \hrulefill 
		%\put(-267,5.5){%
        \put(-193,5.5){
			\makebox(0,0)[l]{%
                \small Faculty of Information and Communication Technology
				%\small Faculty of Information and Communication Technology
			}
		}
	}
	\fancyfoot[LE,RO]{							% Footers
		\thepage
	}
	\renewcommand{\headrulewidth}{0pt}			% Header line width
	\renewcommand{\footrulewidth}{0.2pt}		% Footer line width
}

% Default style for the List of Figures
\fancypagestyle{ListofFiguresStyle}{
	\fancyhf{}									% clear footers and headers
	\fancyhead[RO]{								% Right, odd header
		\hrulefill \hspace{16pt} \large List of Figures
        %\hrulefill \hspace{16pt} \large List of Figures
		\put(-471.5,5.5){%
			\makebox(0,0)[l]{%
                \small Wrocław University of Science and Technology
				%\small Wrocław University of Science and Technology
			}
		}
	}
	\fancyhead[LE]{								% Left, even header
		\large List of Figures \hspace{16pt} \hrulefill 
        %\large List of Figures \hspace{16pt} \hrulefill 
		%\put(-267,5.5){%
        \put(-193,5.5){
			\makebox(0,0)[l]{%
                \small Faculty of Information and Communication Technology
				%\small Faculty of Information and Communication Technology
			}
		}
	}
	\fancyfoot[LE,RO]{							% Footers
		\thepage
	}
	\renewcommand{\headrulewidth}{0pt}			% Header line width
	\renewcommand{\footrulewidth}{0.2pt}		% Footer line width
}

% Default style for the List of Tables
\fancypagestyle{ListofTablesStyle}{
	\fancyhf{}									% clear footers and headers
	\fancyhead[RO]{								% Right, odd header
		\hrulefill \hspace{16pt} \large List of Tables
        % \hrulefill \hspace{16pt} \large List of Tables
		\put(-471.5,5.5){%
			\makebox(0,0)[l]{%
                \small Wrocław University of Science and Technology
				%\small Wrocław University of Science and Technology
			}
		}
	}
	\fancyhead[LE]{								% Left, even header
		\large List of Tables \hspace{16pt} \hrulefill 
        % \large List of Tables \hspace{16pt} \hrulefill 
		%\put(-267,5.5){%
        \put(-193,5.5){
			\makebox(0,0)[l]{%
                \small Faculty of Information and Communication Technology
				%\small Faculty of Information and Communication Technology
			}
		}
	}
	\fancyfoot[LE,RO]{							% Footers
		\thepage
	}
	\renewcommand{\headrulewidth}{0pt}			% Header line width
	\renewcommand{\footrulewidth}{0.2pt}		% Footer line width
}


% Default style for the bibliography
\fancypagestyle{bibliographyStyle}{
	\fancyhf{}									% clear footers and headers
	\fancyhead[RO]{								% Right, odd header
		\hrulefill \hspace{16pt} \large Bibliography
        % \hrulefill \hspace{16pt} \large Bibliography
		\put(-471.5,5.5){%
			\makebox(0,0)[l]{%
				\small Wrocław University of Science and Technology
				%\small Wrocław University of Science and Technology
			}
		}
	}
	\fancyhead[LE]{								% Left, even header
		\large Bibliography \hspace{16pt} \hrulefill 
        % \large Bibliography \hspace{16pt} \hrulefill 
		%\put(-267,5.5){%
        \put(-193,5.5){
			\makebox(0,0)[l]{%
				\small Faculty of Information and Communication Technology
				%\small Faculty of Information and Communication Technology
			}
		}
	}
	\fancyfoot[LE,RO]{							% Footers
		\thepage
	}
	\renewcommand{\headrulewidth}{0pt}			% Header line width
	\renewcommand{\footrulewidth}{0.2pt}		% Footer line width
}

% Default style for appendices
\fancypagestyle{appendixStyle}{
	\fancyhf{}									% clear footers and headers
	\fancyhead[RO]{								% Right, odd header
		\hrulefill \hspace{16pt} \large Appendix \thechapter
        % \hrulefill \hspace{16pt} \large Appendix \thechapter
		\put(-472.1, 12.1){%
			\makebox(0,0)[l]{%
				\includegraphics[width=0.05\textwidth]{resources/pwr-logo}
			}
		}
		\put(-443,5.5){%
			\makebox(0,0)[l]{%
				\small Wrocław University of Science and Technology
				%\small Wrocław University of Science and Technology
			}
		}
	}
	\fancyhead[LE]{								% Left, even header
		\large Appendix \thechapter \hspace{16pt} \hrulefill 
        %\large Appendix \thechapter \hspace{16pt} \hrulefill 
		\put(-22, 12.1){%
			\makebox(0,0)[l]{%
				\includegraphics[width=0.05\textwidth]{wiz-logo}
			}
		}%
		%\put(-220,5.5){%
		\put(-146,5.5){
			\makebox(0,0)[l]{%
                \small Faculty of Information and Communication Technology
				%\small Faculty of Information and Communication Technology
			}
		}
	}
	\fancyfoot[LE,RO]{							% Footers
		\thepage
	}
	\renewcommand{\headrulewidth}{0pt}			% Header line width
	\renewcommand{\footrulewidth}{0.2pt}		% Footer line width
}


% Separate style for pages beginning a chapter/table of contents etc. (formatted as "plain" by default)
\fancypagestyle{chapterBeginStyle}{
	\fancyhf{}%
	\fancyfoot[LE,RO]{
		\thepage
	}
	\renewcommand{\headrulewidth}{0pt}
	\renewcommand{\footrulewidth}{0.2pt}
}

% Style for the remaining table of contents pages
\fancypagestyle{tableOfContentStyle}{
	\fancyhf{}%
	\fancyfoot[LE,RO]{
		\thepage
	}
	\renewcommand{\headrulewidth}{0pt}
	\renewcommand{\footrulewidth}{0.2pt}
}

% Formatting of chapter and/or section titles
\usepackage{titlesec}


% Formatting of chapter titles
\titleformat{\chapter}[hang]
{\normalfont\Huge\bfseries}
{\thechapter.}
{1ex}
{}
\titlespacing*{\chapter}{0pt}{-35pt}{18pt}

% Formatting of section titles
\titleformat{\section}[hang]
{\normalfont\Large\bfseries}											
{\thesection.}
{1ex}
{} 

\titleformat{\subsection}[hang]
{\normalfont\large\bfseries}											
{\thesubsection.}
{1ex}
{} 
% formatting of elements before the modified title


% ***************************************************************************************************
% Links
% ***************************************************************************************************

% Enables insertion of hyperlinks into the document
\usepackage{hyperref}							% Activates links

\hypersetup{
	colorlinks	=	true,					% Colours text instead of creating frames.
	linkcolor	=	blue,					% Colours: references,
    citecolor	=	blue,					% citations,
	urlcolor	=	blue					% hyperlinks.
}

% The same font(s) as the rest of the document will be used to create hyperlinks
\urlstyle{same}


% ***************************************************************************************************
% Links
% ***************************************************************************************************

% Enables defining a custom list style
\usepackage{enumitem}

% New numbered list with three levels
\newlist{myitemize}{itemize}{3}

% Definition of the first-level marker appearance
\setlist[myitemize,1]{
	label		=	\textbullet,
	leftmargin	=	4mm}

% Definition of the second-level marker appearance
\setlist[myitemize,2]{
	label		=	$\diamond$,
	leftmargin	=	8mm}

% Definition of the third-level marker appearance
\setlist[myitemize,3]{
	label		=	$\diamond$,
	leftmargin	=	12mm
}

% ***************************************************************************************************
% Other packages
% ***************************************************************************************************

% Including figures
\usepackage{graphicx}

% Figures and captions
\usepackage{caption}
\usepackage{subcaption}

% Enables creating footnotes inside environments
\usepackage{footnote}

% Enables creating directory structures
\usepackage{dirtree}

% Stretching table cells across multiple rows
\usepackage{multirow}

% Precise calculation of the width/height of any fragment generated by LaTeX
\usepackage{calc}

% ***************************************************************************************************
% Mathematical abbreviations
% ***************************************************************************************************

% Short symbol for real numbers
\newcommand{\RR}{\mathbb{R}}

% Short symbol for natural numbers
\newcommand{\NN}{\mathbb{N}}

% Short symbol for rational numbers
\newcommand{\QQ}{\mathbb{Q}}

% Short symbol for integers
\newcommand{\ZZ}{\mathbb{Z}}

% Short symbol for logical implication
\newcommand{\IMP}{\rightarrow}

% Short symbol for logical equivalence
\newcommand{\IFF}{\leftrightarrow}

% ***************************************************************************************************
% Environments
% ***************************************************************************************************

% Environment for theorems
\newtheorem{theorem}{Theorem}[chapter]

% Environment for lemmas
\newtheorem{lemma}{Lemma}[chapter]

% Environment for examples
\newtheorem{example}{Example}[chapter]

% Environment for corollaries
\newtheorem{corollary}{Corollary}[chapter]

% Environment for definitions
\newtheorem{definition}{Definition}[chapter]

% Environment for proofs
\newenvironment{proof}{
	\par\noindent \textbf{Proof.}
}{
\begin{flushright}
	\vspace*{-6mm}\mbox{$\blacklozenge$}
\end{flushright}
}

% Environment for remarks
\newenvironment{remark}{
	\bigskip \par\noindent \small \textbf{Remark.}
}{
\begin{small}
	\vspace*{4mm}
\end{small}
}

% additional helpers; of course not all are required
\usepackage{psfrag}
\usepackage{amsfonts}
\usepackage{supertabular}
\usepackage{array}
\usepackage{tabularx}
\usepackage{hhline}
\usepackage{minted}
\usepackage{xurl}
\usepackage{microtype}
\usepackage{booktabs} % for professional tables
\usepackage{makecell}
\usepackage{rotating}
\usepackage{multicol}
\usepackage{cuted}
\usepackage{colortbl}
\usepackage{adjustbox}
\usepackage{color,soul}
% \usepackage{subfigure}
\usepackage{pdfpages}

\let\cleardoublepage\clearpage

\usepackage{pifont}
\newcommand{\cmark}{\ding{51}}
\newcommand{\xmark}{\ding{55}}
\newcommand{\bftab}{\fontseries{b}\selectfont}

\newcolumntype{P}[1]{>{\raggedright\arraybackslash\noindent}p{#1}}


\newcolumntype{R}[1]{>{\raggedleft\arraybackslash}p{#1}}
\newcolumntype{L}[1]{>{\raggedright\arraybackslash}p{#1}}
\newcolumntype{C}[1]{>{\centering\arraybackslash}m{#1}}

\usepackage{caption}
\usepackage[nocompress]{cite}
\usepackage{url}
\usepackage{color,soul}
\usepackage{svg}
\usepackage{tabto}
\usepackage{wrapfig}

% formatting of the first pages of chapters - helpers
\newcommand{\resetformatting}{ 
\fancypagestyle{plain}{
    	\fancyhf{}%
    	\fancyfoot[LE,RO]{
    		\thepage
    	}
    	\renewcommand{\headrulewidth}{0pt}
    	\renewcommand{\footrulewidth}{0.2pt}
    } }
\newcommand{\doublepage}{ 
\newpage
\thispagestyle{empty}
\cleardoublepage}

\newcommand\Chapter[1]{ 
\chapter{#1}
\thispagestyle{chapterBeginStyle}}

%***************************************************************************************************
% Custom commands
% ***************************************************************************************************
\newcommand{\kgMargin} [1]{\marginpar{\color{magenta} {\bf PP}: \small{#1}}}
\newcommand{\kg}[1]{{\color{magenta} {\bf PP}: \small{#1}}}

\newcommand{\rrMargin} [1]{\marginpar{\color{magenta} {\bf RR}: \small{#1}}}
\newcommand{\rr}[1]{{\color{blue} {\bf RR}: \small{#1}}}

% Add the remaining comments from the other people - let everyone choose a colour

% Line numbering in the draft
% \usepackage{lineno}
% \renewcommand\linenumberfont{\normalfont\scriptsize}
% \linenumbers
